\documentclass[11pt]{article}
\usepackage[a4paper, top=18mm, bottom=18mm, left=20mm, right=20mm]{geometry}
\usepackage[T2A]{fontenc}
\usepackage[utf8]{inputenc}
\usepackage[ukrainian]{babel}
\usepackage{graphicx}
\usepackage[export]{adjustbox}
\usepackage{amsmath}
\usepackage{amssymb}
\graphicspath{{/app/Quiz/Scripts/2023/ProgAll/15BinTree/}{/app/Quiz/Scripts/2021/Mod2021_3/BinTree/}{/app/Quiz/Scripts/2020/Mod2020_4/BinTree/}}
\setlength{\parindent}{0pt}
\setlength{\parskip}{6pt}
\pagestyle{empty}
\begin{document}
%begin
{\large\textbf{Контрольна робота}}

\textbf{Варіант \No\,3}\hfill \textit{ПІБ:}\,\underline{\hspace{60mm}}
\vspace{4mm}

%beginitem
1. Що буде виведено за виконання фрагменту коду?

%code
print('R', end=' ')
j=1
while j<=9:
    j+=2
print(j)
%code

\vspace{5mm}
%enditem

%beginitem
2. 
5.0<="9.0j"\vspace{5mm}
%enditem

%beginitem
3. 
try:
    res = 7//7/7*7
    if isinstance(res, bool):
        print(f'bool;{res}')
    elif isinstance(res, int):
        print(f'int;{res}')
    elif isinstance(res, float):
        print(f'float;{res:.2f}')
    else:
        print('???')
except:
    print('Z')
\vspace{5mm}
%enditem

%beginitem
4. 
try:
    res = 3**-1--1
    if isinstance(res, bool):
        print(f'bool;{res}')
    elif isinstance(res, int):
        print(f'int;{res}')
    elif isinstance(res, float):
        print(f'float;{res:.2f}')
    else:
        print('???')
except:
    print('Z')
\vspace{5mm}
%enditem

%beginitem
5. 
try:
    for b in range(-3, -3, 3):
        if b<=4:
            continue
        print(b, end=';');b = 2
        if b<=4:
            break
    else:
        print(b,end=';')
    print(b)
except:
    print('Z')
\vspace{5mm}
%enditem

%beginitem
6. %goodpath:Scripts/2019/Mod2019_1/Lex/03-good.txt
Відмітити все, що є коректним оператором С++:
( 1 ) return 'a'<'c'

( 2 ) print(sep=":","qwe")

( 3 ) x=*2

( 4 ) float(2)/87
\vspace{5mm}
%enditem

%beginitem
7. %goodpath:Scripts/2019/Mod2019_1/Lex/01-good.txt
Відмітити все, що є однією правильною лексемою:
( 1 ) 'c'

( 2 ) 0123

( 3 ) 0b23

( 4 ) xy
\vspace{5mm}
%enditem

%beginitem
8. Що буде виведено за виконання фрагменту коду?

%code
cout << 8 / 8 / 8 / 8;
%code
\vspace{5mm}
%enditem

%beginitem
9. %goodpath:Scripts/2018/Mod2018_1/01-good.txt
Відмітити все, що є однією правильною лексемою:
( 1 ) a2f

( 2 ) 'qwerty'

( 3 ) **

( 4 ) true
\vspace{5mm}
%enditem

%beginitem
10. Що буде виведено за виконання фрагменту коду?
%code
#include <iostream>
using namespace std;

int f(int &x, int &y){
    x = 5;
    y-= 9;
    return y;
}

int main(){
    int a = 6, b = 4;
    b = f(b, b);
    cout << a << ":" << b <<':';
    {
        int a = 2, b = 3;
        cout << ((b<=7) || ((a-=2) <= 7))
             << a << ":" << b << ":";
    }
    {
        inta = 1;
        cout << a << ':';
    }
    cout << a << ":" << b << endl;
    return 0;
}
%code

\vspace{5mm}
%enditem

%end
\newpage
\end{document}

